%% ===================================================================================
%% Capítulo 9 -- Considerações Finais
%% Base: Conclusão do relatório de IC (resultados, limitações, trabalhos futuros),
%% ampliada para o TCC. As conclusões sobre o monólito são escritas pelo Arquiteto
%% somente após o Capítulo 8 estar preenchido.
%% ===================================================================================
\chapter{Considerações Finais}\label{cap:conclusao}

\section{Conclusões}\label{sec:conclusoes}

Este trabalho tratou do desenvolvimento do Vitrine, uma plataforma CMS de código aberto voltada à publicação acadêmica independente. Partiu de um problema concreto --- a dependência de plataformas proprietárias ou a complexidade de CMS tradicionais --- e de uma base já concebida e prototipada na Iniciação Científica \cite{reis2026ic}, que aponta a soberania do autor sobre a sua produção intelectual como princípio orientador.

A contribuição deste trabalho é dupla. No plano do produto, a consolidação do sistema em um monólito modular sobre Next.js, SQLite e better-auth, descrito por requisitos rastreáveis, modelo de dados, contrato de API, estratégia de cache e decisões de projeto (Capítulos~\ref{cap:requisitos} e~\ref{cap:arquitetura}). No plano do processo, a proposta e a aplicação de um método de \en{Documentation-Driven Development} em que uma documentação viva, versionada junto ao código e com um único redator, coordena agentes especializados (Capítulo~\ref{cap:metodologia}).

\vitrinependencia{Arquiteto}{Após o preenchimento do Capítulo 8, reescrever estas conclusões retomando o objetivo geral e cada objetivo específico (OE1 a OE8) com base nas evidências, e responder às perguntas de pesquisa P1 e P2 da Seção 1.3.}

\section{Contribuições}\label{sec:contribuicoes-finais}

\begin{enumerate}
  \item Um artefato de \en{software} de código aberto para publicação científica autônoma, com edição Markdown e fórmulas \LaTeX{}, personalização visual dinâmica e entrega por ISR;
  \item uma especificação documentada e rastreável do sistema, reutilizável por terceiros;
  \item um método de desenvolvimento por documentação e agentes, com regras de governança, fluxo de Git e mural de comunicação, descrito de forma replicável;
  \item um relato de verificação e validação, com as limitações e as ameaças à validade.
\end{enumerate}

\section{Limitações}\label{sec:limitacoes-finais}

Além das limitações registradas na Seção~\ref{sec:res-ameacas}, o relatório de IC já apontava que a modelagem da interface se restringira a protótipos de baixa fidelidade, em prol da funcionalidade e da usabilidade, e que a autenticação poderia ser expandida para múltiplos níveis de permissão em cenários colaborativos \cite{reis2026ic}. Ambos os pontos permanecem como oportunidades de evolução.

\section{Trabalhos Futuros}\label{sec:trabalhos-futuros}

Do relatório de IC \cite{reis2026ic} permanecem como propostas de continuidade:

\begin{itemize}
  \item produzir documentação técnica detalhada e tutoriais em vídeo que demonstrem o processo de clonagem do repositório e de implantação --- considerando a ressalva sobre plataformas \emph{serverless} e a persistência do SQLite (\decisao{D-04});
  \item lançar o site oficial do projeto, denominado Vitrine, que servirá como repositório central para tutoriais, documentação e instâncias de demonstração;
  \item explorar a integração com repositórios de dados abertos para a citação automática de publicações;
  \item expandir a biblioteca de componentes visuais, permitindo maior customização da interface;
  \item implementar a renderização nativa de gráficos de funções de uma variável (gráficos 2D).
\end{itemize}

Da consolidação do monólito surgem novas linhas de evolução, já mapeadas como decisões em aberto (Apêndice~\ref{ap:decisoes}): papéis e permissões com o \en{plugin} \texttt{admin} do better-auth (\decisao{D-13}); estados de publicação e histórico de versões (\decisao{D-15}); envio e gerenciamento de mídia (\decisao{D-10}); busca textual no conteúdo; internacionalização (\req{RNF12}); exportação e importação do conteúdo em Markdown, para reforçar a soberania dos dados; e temas pré-definidos.

\section{Considerações Finais}

A experiência de desenvolver o Vitrine reafirma a premissa de que sistemas de informação devem atuar como facilitadores da disseminação de conhecimento, e não como obstáculos técnicos \cite{rainer2019intro}. O projeto entrega uma base sólida e moderna que elimina o trabalho repetitivo de configuração estrutural, permitindo que o foco permaneça na produção de conteúdo qualificado. O Vitrine busca ser não apenas uma ferramenta técnica, mas um instrumento em prol da autonomia digital e da democratização do saber científico na rede.
